\chapter{2006年陕西卷（文）}
\section{单选题}
\begin{problem}
已知集合 \(P = \{x \in \N \mid 1 \leqslant x \leqslant 10\}\)，集合 \(Q = \{x \in \R \mid x^2 + x - 6 = 0\}\)，则 \(P \cap Q\) 等于（\quad）

\choices
    {\(\{2\}\)}
    {\(\{3\}\)}
    {\(\{-2, 3\}\)}
    {\(\{-3, 2\}\)}
\end{problem}

\begin{problem}
函数 \( f(x) = \frac{1}{1 + x^2} (x \in \R) \) 的值域是（\quad）

\choices
    {\((0,1)\)}
    {\((0,1]\)}
    {\([0,1)\)}
    {\([0,1]\)}
\end{problem}

\begin{problem}
已知等差数列 \(\{a_{n}\}\) 中，\(a_{1} + a_{8} = 8\)，则该数列前 9 项和 \(S_{9}\) 等于（\quad）

\choices
    {18}
    {27}
    {36}
    {45}
\end{problem}

\begin{problem}
设函数 \( f(x) = \log_a(x + b) (a > 0, a \neq 1) \) 的图象过点 \((0,0)\)，其反函数的图象过点 \((1,2)\)，则 \( a + b \) 等于（\quad）

\choices
    {6}
    {5}
    {4}
    {3}
\end{problem}

\begin{problem}
设直线过点 \((0, a)\)，其斜率为 1，且与圆 \(x^{2} + y^{2} = 2\) 相切，则 \(a\) 的值为（\quad）

\choices
    {\(\pm \sqrt{2}\)}
    {\(\pm 2\)}
    {\(\pm 2\sqrt{2}\)}
    {\(\pm 4\)}
\end{problem}

\begin{problem}
“\(\alpha\)，\(\beta\)，\(\gamma\) 成等差数列”是“等式 \(\sin (\alpha + \gamma) = \sin 2\beta\) 成立”的（\quad）

\choices
    {充分而不必要条件}
    {必要而不充分条件}
    {充分必要条件}
    {既不充分也不必要条件}
\end{problem}

\begin{problem}
设 \(x\)，\(y\) 为正数，则 \((x + y)\left(\frac{1}{x} + \frac{4}{y}\right)\) 的最小值为（\quad）

\choices
    {6}
    {9}
    {12}
    {15}
\end{problem}

\begin{problem}
已知非零向量 \(\overrightarrow{AB}\) 与 \(\overrightarrow{AC}\) 满足 \(\left(\frac{\overrightarrow{AB}}{|\overrightarrow{AB}|} + \frac{\overrightarrow{AC}}{|\overrightarrow{AC}|}\right)\cdot \overrightarrow{BC} = 0\) 且 \(\frac{\overrightarrow{AB}}{|\overrightarrow{AB}|}\cdot \frac{\overrightarrow{AC}}{|\overrightarrow{AC}|} = \frac{1}{2}\)，则 \(\triangle ABC\) 为（\quad）

\choices
    {三边均不相等的三角形}
    {直角三角形}
    {等腰非等边三角形}
    {等边三角形}
\end{problem}

\begin{problem}
已知函数 \( f(x) = ax^2 + 2ax + 4 (a > 0) \). 若 \(x_1 < x_2\)，\(x_1 + x_2 = 0\)，则（\quad）

\choices
    {\(f(x_{1}) < f(x_{2})\)}
    {\(f(x_{1}) = f(x_{2})\)}
    {\(f(x_{1}) > f(x_{2})\)}
    {\(f(x_{1})\) 与 \(f(x_{2})\) 的大小不能确定}
\end{problem}

\begin{problem}
已知双曲线 \(\frac{x^2}{a^2} - \frac{y^2}{2} = 1 (a > \sqrt{2})\) 的两条渐近线的夹角为 \(\frac{\pi}{3}\)，则双曲线的离心率为（\quad）

\choices
    {2}
    {\(\sqrt{3}\)}
    {\(\frac{2\sqrt{6}}{3}\)}
    {\(\frac{2\sqrt{3}}{3}\)}
\end{problem}

\begin{problem}
已知平面 \(\alpha\) 外不共线的三点 \(A\)，\(B\)，\(C\) 到 \(\alpha\) 的距离都相等. 则正确的结论是（\quad）

\choices
    {平面 \(ABC\) 必平行于 \(\alpha\)}
    {平面 \(ABC\) 必与 \(\alpha\) 相交}
    {平面 \(ABC\) 必不垂直于 \(\alpha\)}
    {存在 \(\triangle ABC\) 的一条中位线平行于 \(\alpha\) 或在 \(\alpha\) 内}
\end{problem}

\begin{problem}
为确保信息安全，信息需加密传播，发送方由明文 \(\rightarrow\) 密文（加密），接收方由密文 \(\rightarrow\) 明文（解密）. 已知加密规则为：明文 \(a\)，\(b\)，\(c\)，\(d\) 对应密文 \(a + 2b\)，\(2b + c\)，\(2c + 3d\)，\(4d\). 例如，明文 1， 2， 3， 4 对应密文 5， 7， 18， 16. 当接收方收到密文 14， 9， 23， 28 时，则解密得到的明文为（\quad）

\choices
    {1, 6, 4, 7}
    {4, 6, 1, 7}
    {7, 6, 1, 4}
    {6, 4, 1, 7}
\end{problem}

\section{填空题}
\begin{problem}
\(\cos 43^{\circ}\cos 77^{\circ} + \sin 43^{\circ}\cos 167^{\circ}\) 的值为
\end{problem}

\begin{problem}
\(\left(2x - \frac{1}{\sqrt{x}}\right)^6\) 展开式中的常数项为\fillinblank{}. （用数字作答）
\end{problem}

\begin{problem}
某校从 8 名教师中选派 4 名教师同时去 4 个边远地区支教（每地 1 人），其中甲和乙不同去，则不同的选派方案共有\fillinblank{} 种.
\end{problem}

\begin{problem}
水平桌面 \(\alpha\) 上放有 4 个半径为 \(2R\) 的球. 且相邻的球都相切（球心的连线构成正方形）. 在这 4 个球的上面放 1 个半径为 \(R\) 的小球，它和下面的 4 个球恰好都相切，则小球的球心到水平桌面 \(\alpha\) 的距离是\fillinblank{}.
\end{problem}

\section{解答题}
\begin{problem}
甲，乙，丙 3 人投篮，投进的概率分别是 \(\frac{2}{5}\)，\(\frac{1}{2}\)，\(\frac{3}{5}\)，现 3 人各投篮 1 次，求：

\begin{enumerate}
\item 3 人都投进的概率；
\item 3 人中恰有 2 人投进的概率.
\end{enumerate}
\end{problem}

\begin{problem}
已知函数 \( f(x) = \sqrt{3}\sin \left(2x - \frac{\pi}{6}\right) + 2\sin^2\left(x - \frac{\pi}{12}\right) (x \in \R) \).

\begin{enumerate}
\item 求函数 \( f(x) \) 的最小正周期；
\item 求使函数 \( f(x) \) 取得最大值的 \(x\) 的集合.
\end{enumerate}
\end{problem}

\begin{problem}
如图，\(\alpha \perp \beta\)，\(\alpha \cap \beta = l\)，\(A \in \alpha\)，\(B \in \beta\)，点 \(A\) 在直线 \(l\) 上的射影为 \(A_1\)，点 \(B\) 在 \(l\) 上的射影为 \(B_1\). 已知 \(AB = 2\)，\(AA_1 = 1\)，\(BB_1 = \sqrt{2}\). 求：

\begin{enumerate}
\item 直线 \(AB\) 分别与平面 \(\alpha\)，\(\beta\) 所成角的大小；
\item 二面角 \(A_{1} - AB - B_{1}\) 的大小.
\end{enumerate}

\bitmapfigure[max width=0.38\linewidth,max totalheight=4.8cm]{img/2006/shaanxi_liberal/q19_fig1.png}

\end{problem}

\begin{problem}
已知正项数列 \(\{a_{n}\}\)，其前 \(n\) 项和 \(S_{n}\) 满足 \(10S_{n} = a_{n}^{2} + 5a_{n} + 6\)，且 \(a_{1}\)，\(a_{3}\)，\(a_{15}\) 成等比数列，求数列 \(\{a_{n}\}\) 的通项 \(a_{n}\).
\end{problem}

\begin{problem}
如图，三定点 \(A(2,1)\)，\(B(0,-1)\)，\(C(-2,1)\)，三动点 \(D\)，\(E\)，\(M\) 满足 \(\overrightarrow{AD} = t\overrightarrow{AB}\)，\(\overrightarrow{BE} = t\overrightarrow{BC}\)，\(\overrightarrow{DM} = t\overrightarrow{DE}\)，\(t \in [0,1]\). 求：

\begin{enumerate}
\item 求动直线 \(DE\) 斜率的变化范围；
\item 求动点 \(M\) 的轨迹方程.
\end{enumerate}

\bitmapfigure[max width=0.40\linewidth,max totalheight=4.8cm]{img/2006/shaanxi_liberal/q21_fig1.png}

\end{problem}

\begin{problem}
设函数 \( f(x) = kx^{3} - 3x^{2} + 1 (k \geqslant 0) \).

\begin{enumerate}
\item 求函数 \( f(x) \) 的单调区间；
\item 若函数 \( f(x) \) 的极小值大于 0，求 \(k\) 的取值范围.
\end{enumerate}
\end{problem}
